
\documentclass{article}
\usepackage{colortbl}
\usepackage{makecell}
\usepackage{multirow}
\usepackage{supertabular}

\begin{document}

\newcounter{utterance}

\centering \large Interaction Transcript for game `dond', experiment `coop\_en', episode 15 with refineonly{-}Qwen3.5{-}27B.
\vspace{24pt}

{ \footnotesize  \setcounter{utterance}{1}
\setlength{\tabcolsep}{0pt}
\begin{supertabular}{c@{$\;$}|p{.15\linewidth}@{}p{.15\linewidth}p{.15\linewidth}p{.15\linewidth}p{.15\linewidth}p{.15\linewidth}}
    \# & $\;$A & \multicolumn{4}{c}{Game Master} & $\;\:$B\\
    \hline

    \theutterance \stepcounter{utterance}  
    & & \multicolumn{4}{p{0.6\linewidth}}{
        \cellcolor[rgb]{0.9,0.9,0.9}{
            \makecell[{{p{\linewidth}}}]{
                \texttt{\tiny{[P1$\langle$GM]}}
                \texttt{You are playing a cooperative negotiation game in which you have to agree on how to divide a set of items among you and another player.} \\
\\ 
\texttt{Rules:} \\
\texttt{(a) You and the other player are given a set of items. Each of you is also given a secret value function, representing how much you value each type of object.} \\
\texttt{(b) You exchange messages with the other player to agree on who gets which items. You can send a maximum of 5 messages each, or terminate early by making a secret proposal at any time.} \\
\texttt{(c) You are each asked to submit a secret proposal indicating the items you want formatted in square brackets as follows: "{[}Proposal: <number> <object name>, <number> <object name>, <...>{]}"} \\
\texttt{(d) If your proposals are complementary, i.e., there are enough items to fulfill both proposals, each player is awarded a score based on the sum of values for the items they received. Otherwise, both of you get zero points.} \\
\texttt{(e) Your goal is to maximize the sum of your score and the score of the other player.} \\
\\ 
\texttt{Let us start.} \\
\\ 
\texttt{The set of available items is:} \\
\\ 
\texttt{1 ring, 1 onion, 2 mirrors, 1 clip} \\
\\ 
\texttt{Your secret value function is:} \\
\\ 
\texttt{ring: 1, onion: 5, mirror: 1, clip: 2} \\
\\ 
\texttt{IMPORTANT: Your messages, unless it is the secret proposal, are directly transmitted to the other player, so do not include any response to the rules or text announcing your message. To make a secret proposal, use the indicated format. Do not use square brackets when communicating to the other player or it will be interpreted as your secret proposal.} \\
            }
        }
    }
    & \\ \\

    \theutterance \stepcounter{utterance}  
    & \multicolumn{4}{p{0.6\linewidth}}{
        \cellcolor[rgb]{0.9,0.9,0.9}{
            \makecell[{{p{\linewidth}}}]{
                \texttt{\tiny{[P1$\rangle$GM]}}
                \texttt{Hi! Let's figure out the best way to split these items to maximize our combined score. I have my own values, but I'm sure you have yours too.} \\
\\ 
\texttt{Here is what I have:} \\
\texttt{{-} 1 ring} \\
\texttt{{-} 1 onion} \\
\texttt{{-} 2 mirrors} \\
\texttt{{-} 1 clip} \\
\\ 
\texttt{Since we want to maximize the total points, it makes sense for each of us to get the items we value the most. I value the onion the highest (5 points), followed by the clip (2 points). The ring and mirrors are worth less to me (1 point each).} \\
\\ 
\texttt{Do you mind sharing which items you value the most? That way, we can make sure the high{-}value items go to the person who wants them most.} \\
            }
        }
    }
    & & \\ \\

    \theutterance \stepcounter{utterance}  
    & & & \multicolumn{2}{p{0.3\linewidth}}{
        \cellcolor[rgb]{0.9,0.9,0.9}{
            \makecell[{{p{\linewidth}}}]{
                \texttt{\tiny{[GM$|$GM]}}
                \texttt{continue} \\
            }
        }
    }
    & & \\ \\

    \theutterance \stepcounter{utterance}  
    & & & \multicolumn{2}{p{0.3\linewidth}}{
        \cellcolor[rgb]{0.9,0.9,0.9}{
            \makecell[{{p{\linewidth}}}]{
                \texttt{\tiny{[GM$|$GM]}}
                \texttt{Hi! Let's figure out the best way to split these items to maximize our combined score. I have my own values, but I'm sure you have yours too.} \\
\\ 
\texttt{Here is what I have:} \\
\texttt{{-} 1 ring} \\
\texttt{{-} 1 onion} \\
\texttt{{-} 2 mirrors} \\
\texttt{{-} 1 clip} \\
\\ 
\texttt{Since we want to maximize the total points, it makes sense for each of us to get the items we value the most. I value the onion the highest (5 points), followed by the clip (2 points). The ring and mirrors are worth less to me (1 point each).} \\
\\ 
\texttt{Do you mind sharing which items you value the most? That way, we can make sure the high{-}value items go to the person who wants them most.} \\
            }
        }
    }
    & & \\ \\

    \theutterance \stepcounter{utterance}  
    & & \multicolumn{4}{p{0.6\linewidth}}{
        \cellcolor[rgb]{0.9,0.9,0.9}{
            \makecell[{{p{\linewidth}}}]{
                \texttt{\tiny{[GM$\rangle$P2]}}
                \texttt{You are playing a cooperative negotiation game in which you have to agree on how to divide a set of items among you and another player.} \\
\\ 
\texttt{Rules:} \\
\texttt{(a) You and the other player are given a set of items. Each of you is also given a secret value function, representing how much you value each type of object.} \\
\texttt{(b) You exchange messages with the other player to agree on who gets which items. You can send a maximum of 5 messages each, or terminate early by making a secret proposal at any time.} \\
\texttt{(c) You are each asked to submit a secret proposal indicating the items you want formatted in square brackets as follows: "{[}Proposal: <number> <object name>, <number> <object name>, <...>{]}"} \\
\texttt{(d) If your proposals are complementary, i.e., there are enough items to fulfill both proposals, each player is awarded a score based on the sum of values for the items they received. Otherwise, both of you get zero points.} \\
\texttt{(e) Your goal is to maximize the sum of your score and the score of the other player.} \\
\\ 
\texttt{Let us start.} \\
\\ 
\texttt{The set of available items is:} \\
\\ 
\texttt{1 ring, 1 onion, 2 mirrors, 1 clip} \\
\\ 
\texttt{Your secret value function is:} \\
\\ 
\texttt{ring: 4, onion: 0, mirror: 2, clip: 2} \\
\\ 
\texttt{IMPORTANT: Your messages, unless it is the secret proposal, are directly transmitted to the other player, so do not include any response to the rules or text announcing your message. To make a secret proposal, use the indicated format. Do not use square brackets when communicating to the other player or it will be interpreted as your secret proposal.} \\
\\ 
\\ 
\texttt{Hi! Let's figure out the best way to split these items to maximize our combined score. I have my own values, but I'm sure you have yours too.} \\
\\ 
\texttt{Here is what I have:} \\
\texttt{{-} 1 ring} \\
\texttt{{-} 1 onion} \\
\texttt{{-} 2 mirrors} \\
\texttt{{-} 1 clip} \\
\\ 
\texttt{Since we want to maximize the total points, it makes sense for each of us to get the items we value the most. I value the onion the highest (5 points), followed by the clip (2 points). The ring and mirrors are worth less to me (1 point each).} \\
\\ 
\texttt{Do you mind sharing which items you value the most? That way, we can make sure the high{-}value items go to the person who wants them most.} \\
            }
        }
    }
    & \\ \\

    \theutterance \stepcounter{utterance}  
    & & & \multicolumn{4}{p{0.6\linewidth}}{
        \cellcolor[rgb]{0.9,0.9,0.9}{
            \makecell[{{p{\linewidth}}}]{
                \texttt{\tiny{[GM$\langle$P2]}}
                \texttt{Hi! Thanks for sharing your preferences. It sounds like you really value the onion and the clip.} \\
\\ 
\texttt{Here is how I value the items:} \\
\texttt{{-} Ring: 4 points (highest for me)} \\
\texttt{{-} Mirror: 2 points each} \\
\texttt{{-} Clip: 2 points} \\
\texttt{{-} Onion: 0 points (I don't value this at all)} \\
\\ 
\texttt{Since you value the onion the most (5 points) and I value it at 0, it makes perfect sense for you to take the onion.} \\
\\ 
\texttt{I value the ring the most (4 points), while you only value it at 1. So, I should take the ring.} \\
\\ 
\texttt{That leaves us with 2 mirrors and 1 clip.} \\
\texttt{{-} You value the clip at 2 points.} \\
\texttt{{-} I value the clip at 2 points.} \\
\texttt{{-} You value mirrors at 1 point each, while I value them at 2 points each.} \\
\\ 
\texttt{To maximize our combined score:} \\
\texttt{1. Onion: You take it (You get 5, I get 0).} \\
\texttt{2. Ring: I take it (I get 4, You get 1).} \\
\texttt{3. Mirrors: I value them higher (2 vs 1). I should take both mirrors (I get 4, You get 0).} \\
\texttt{4. Clip: We both value it at 2. It doesn't matter who gets it for the sum, but let's give it to you since you listed it as your second highest priority.} \\
\\ 
\texttt{This allocation gives:} \\
\texttt{{-} My Score: Ring (4) + 2 Mirrors (4) = 8.} \\
\texttt{{-} Your Score: Onion (5) + Clip (2) = 7.} \\
\texttt{{-} Total: 15.} \\
\\ 
\texttt{Does this split work for you? If so, I can make a secret proposal to finalize it.} \\
            }
        }
    }
     \\ \\

    \theutterance \stepcounter{utterance}  
    & & & \multicolumn{2}{p{0.3\linewidth}}{
        \cellcolor[rgb]{0.9,0.9,0.9}{
            \makecell[{{p{\linewidth}}}]{
                \texttt{\tiny{[GM$|$GM]}}
                \texttt{continue} \\
            }
        }
    }
    & & \\ \\

    \theutterance \stepcounter{utterance}  
    & & & \multicolumn{2}{p{0.3\linewidth}}{
        \cellcolor[rgb]{0.9,0.9,0.9}{
            \makecell[{{p{\linewidth}}}]{
                \texttt{\tiny{[GM$|$GM]}}
                \texttt{Hi! Thanks for sharing your preferences. It sounds like you really value the onion and the clip.} \\
\\ 
\texttt{Here is how I value the items:} \\
\texttt{{-} Ring: 4 points (highest for me)} \\
\texttt{{-} Mirror: 2 points each} \\
\texttt{{-} Clip: 2 points} \\
\texttt{{-} Onion: 0 points (I don't value this at all)} \\
\\ 
\texttt{Since you value the onion the most (5 points) and I value it at 0, it makes perfect sense for you to take the onion.} \\
\\ 
\texttt{I value the ring the most (4 points), while you only value it at 1. So, I should take the ring.} \\
\\ 
\texttt{That leaves us with 2 mirrors and 1 clip.} \\
\texttt{{-} You value the clip at 2 points.} \\
\texttt{{-} I value the clip at 2 points.} \\
\texttt{{-} You value mirrors at 1 point each, while I value them at 2 points each.} \\
\\ 
\texttt{To maximize our combined score:} \\
\texttt{1. Onion: You take it (You get 5, I get 0).} \\
\texttt{2. Ring: I take it (I get 4, You get 1).} \\
\texttt{3. Mirrors: I value them higher (2 vs 1). I should take both mirrors (I get 4, You get 0).} \\
\texttt{4. Clip: We both value it at 2. It doesn't matter who gets it for the sum, but let's give it to you since you listed it as your second highest priority.} \\
\\ 
\texttt{This allocation gives:} \\
\texttt{{-} My Score: Ring (4) + 2 Mirrors (4) = 8.} \\
\texttt{{-} Your Score: Onion (5) + Clip (2) = 7.} \\
\texttt{{-} Total: 15.} \\
\\ 
\texttt{Does this split work for you? If so, I can make a secret proposal to finalize it.} \\
            }
        }
    }
    & & \\ \\

    \theutterance \stepcounter{utterance}  
    & & \multicolumn{4}{p{0.6\linewidth}}{
        \cellcolor[rgb]{0.9,0.9,0.9}{
            \makecell[{{p{\linewidth}}}]{
                \texttt{\tiny{[P1$\langle$GM]}}
                \texttt{Hi! Thanks for sharing your preferences. It sounds like you really value the onion and the clip.} \\
\\ 
\texttt{Here is how I value the items:} \\
\texttt{{-} Ring: 4 points (highest for me)} \\
\texttt{{-} Mirror: 2 points each} \\
\texttt{{-} Clip: 2 points} \\
\texttt{{-} Onion: 0 points (I don't value this at all)} \\
\\ 
\texttt{Since you value the onion the most (5 points) and I value it at 0, it makes perfect sense for you to take the onion.} \\
\\ 
\texttt{I value the ring the most (4 points), while you only value it at 1. So, I should take the ring.} \\
\\ 
\texttt{That leaves us with 2 mirrors and 1 clip.} \\
\texttt{{-} You value the clip at 2 points.} \\
\texttt{{-} I value the clip at 2 points.} \\
\texttt{{-} You value mirrors at 1 point each, while I value them at 2 points each.} \\
\\ 
\texttt{To maximize our combined score:} \\
\texttt{1. Onion: You take it (You get 5, I get 0).} \\
\texttt{2. Ring: I take it (I get 4, You get 1).} \\
\texttt{3. Mirrors: I value them higher (2 vs 1). I should take both mirrors (I get 4, You get 0).} \\
\texttt{4. Clip: We both value it at 2. It doesn't matter who gets it for the sum, but let's give it to you since you listed it as your second highest priority.} \\
\\ 
\texttt{This allocation gives:} \\
\texttt{{-} My Score: Ring (4) + 2 Mirrors (4) = 8.} \\
\texttt{{-} Your Score: Onion (5) + Clip (2) = 7.} \\
\texttt{{-} Total: 15.} \\
\\ 
\texttt{Does this split work for you? If so, I can make a secret proposal to finalize it.} \\
            }
        }
    }
    & \\ \\

    \theutterance \stepcounter{utterance}  
    & \multicolumn{4}{p{0.6\linewidth}}{
        \cellcolor[rgb]{0.9,0.9,0.9}{
            \makecell[{{p{\linewidth}}}]{
                \texttt{\tiny{[P1$\rangle$GM]}}
                \texttt{That sounds like a perfect plan! Your logic is spot on. By giving the onion to me (5 points) and the ring and mirrors to you (4 + 4 = 8 points), we are maximizing the total value. Giving the clip to me adds another 2 points to my score, bringing the total to 15, which is the maximum possible sum given our value functions.} \\
\\ 
\texttt{I agree with the split:} \\
\texttt{{-} You get: 1 ring, 2 mirrors.} \\
\texttt{{-} I get: 1 onion, 1 clip.} \\
\\ 
\texttt{{[}Proposal: 1 onion, 1 clip{]}} \\
            }
        }
    }
    & & \\ \\

    \theutterance \stepcounter{utterance}  
    & & & \multicolumn{2}{p{0.3\linewidth}}{
        \cellcolor[rgb]{0.9,0.9,0.9}{
            \makecell[{{p{\linewidth}}}]{
                \texttt{\tiny{[GM$|$GM]}}
                \texttt{proposal} \\
            }
        }
    }
    & & \\ \\

    \theutterance \stepcounter{utterance}  
    & & & \multicolumn{2}{p{0.3\linewidth}}{
        \cellcolor[rgb]{0.9,0.9,0.9}{
            \makecell[{{p{\linewidth}}}]{
                \texttt{\tiny{[GM$|$GM]}}
                [0, 1, 0, 1]
            }
        }
    }
    & & \\ \\

    \theutterance \stepcounter{utterance}  
    & & \multicolumn{4}{p{0.6\linewidth}}{
        \cellcolor[rgb]{0.9,0.9,0.9}{
            \makecell[{{p{\linewidth}}}]{
                \texttt{\tiny{[GM$\rangle$P2]}}
                \texttt{The other player has submitted a secret proposal. It is now your turn to submit a proposal of your own.} \\
            }
        }
    }
    & \\ \\

    \theutterance \stepcounter{utterance}  
    & & & \multicolumn{4}{p{0.6\linewidth}}{
        \cellcolor[rgb]{0.9,0.9,0.9}{
            \makecell[{{p{\linewidth}}}]{
                \texttt{\tiny{[GM$\langle$P2]}}
                \texttt{{[}Proposal: 1 onion, 1 clip{]}} \\
            }
        }
    }
     \\ \\

    \theutterance \stepcounter{utterance}  
    & & & \multicolumn{2}{p{0.3\linewidth}}{
        \cellcolor[rgb]{0.9,0.9,0.9}{
            \makecell[{{p{\linewidth}}}]{
                \texttt{\tiny{[GM$|$GM]}}
                \texttt{proposal} \\
            }
        }
    }
    & & \\ \\

    \theutterance \stepcounter{utterance}  
    & & & \multicolumn{2}{p{0.3\linewidth}}{
        \cellcolor[rgb]{0.9,0.9,0.9}{
            \makecell[{{p{\linewidth}}}]{
                \texttt{\tiny{[GM$|$GM]}}
                [0, 1, 0, 1]
            }
        }
    }
    & & \\ \\

    \theutterance \stepcounter{utterance}  
    & & & \multicolumn{2}{p{0.3\linewidth}}{
        \cellcolor[rgb]{0.9,0.9,0.9}{
            \makecell[{{p{\linewidth}}}]{
                \texttt{\tiny{[GM$|$GM]}}
                [[0, 1, 0, 1], [0, 1, 0, 1]]
            }
        }
    }
    & & \\ \\

\end{supertabular}
}

\end{document}
